\section{A common reduction at the compression boundary}
\label{sec:three-case-galois-pu}
The next proposition removes a pivotality gap for three quadratic rings.
It supplies the general-complex scope for the common exclusion in
Theorem~\ref{thm:quadratic-center-trace}.

\begin{proposition}\label{prop:three-case-galois-pu}
Every complex fusion categorification of $R7\_4\_106$,
$R7\_4\_2155$, or $R7\_3\_146$ is Galois conjugate to a
pseudo-unitary category.
\end{proposition}
\begin{proof}
Write $g$ for the nontrivial invertible, $a,ga$ and $b,gb$ for its
two free pairs, and $f$ for its fixed noninvertible. The source labels
$(g,a,b,f)$ are $(1,2,3,6)$, $(3,1,5,4)$, and $(1,2,3,6)$,
respectively. Duality is trivial except that $a^*=ga$ in the second ring.
Pass to the signed algebra $T$ of the canonical sphericalization,
at $\varepsilon=-1$. Choose the lift of $g$ with dimension $1$.
Then $g^2=1$, since the other lift of the unit has dimension $-1$.
The same nonzero-dimension argument makes $g$ central and gives
$gf=f$. Choose the free-pair lifts compatibly with multiplication
by $g$. The quotient $T/(g-1)$ has basis $1,a,b,f$ and trivial
duality, including in the second case. It is therefore commutative.
Frobenius reciprocity gives
\begin{align*}
 a^2&=1+Aa+Bb+rf,&ab&=Ba+Cb+wf,\\
 b^2&=1+Ca+Db+zf,&af&=2ra+2wb+lf,\\
 bf&=2wa+2zb+mf,&f^2&=2+2la+2mb+nf.
\end{align*}
The coefficient pairing has weights $1,1,1,2$; the factors two
come from the free pairs. Each coefficient is an integer of the
same parity as its source bound and of absolute value at most that
bound. In order $(A,B,C,D,r,w,z,l,m,n)$ the bounds are
\[
 (0,3,4,8,1,2,4,2,3,4),\qquad
 (3,3,4,4,2,2,2,2,3,2),
\]
for the $\sqrt3$ pair and the $\sqrt6$ case, respectively. Independent
sign changes of $a,b,f$ allow $A>0$ whenever $A$ is odd, $B>0$,
and $r\geq0$. In the first family $A=0$ and $r=1$.

Associativity is equivalent to these six equations:
\begin{align*}
 AC-B^2+BD-C^2+2rz-2w^2+1&=0,\\
 Al+Bm-l^2+nr-2r^2-2w^2+1&=0,\\
 Aw-Br+Bz-Cw-lw+mr&=0,\\
 Bl+Cm-lm+nw-2rw-2wz&=0,\\
 Bw-Cr+Cz-Dw-lz+mw&=0,\\
 Cl+Dm-m^2+nz-2w^2-2z^2+1&=0.
\end{align*}
They have eight and eleven solutions in the respective normalized
domains. To obtain the complete lists, $B\in\{1,3\}$ makes the
first equation determine $D$; if $r\ne0$, the third and second
determine $m,n$. If $r=0$, only their original small domains remain.
The source-bound verifier records all nineteen rows and checks every
multiplication-matrix identity.

Let $\delta$ be the global dimension of the original category.
Its signed dimension character factors through this quotient. If
its values on $a,b,f$ are $u,v,w_0$, then
$\delta=2+2u^2+2v^2+w_0^2$, also in the non-self-dual case because
$a^*=ga$. Thus $\delta$ is an eigenvalue of
\[
                  Q=2I+2M_a^2+2M_b^2+M_f^2.
\]
The induced-unit theorem for the sphericalization gives
$2\delta/(2f_E)=\delta/f_E$ algebraically integral for every
source formal codegree $f_E$ \cite[Theorem~2.13]{OstrikRank3}.
One is $48(7+4\sqrt3)$ in the first family and $48(5+2\sqrt6)$
in the second. Each parenthesized factor is an algebraic unit of
norm one, so multiplication by it proves that $\delta/48$ is
integral. The rational source codegrees alone have least common
multiple $24$ and would not give this stronger divisibility.

For an irreducible factor $p(x)$ of $\det(xI-Q)$, a root divided by
$48$ is integral exactly when the monic polynomial proportional to
$p(48x)$ has integer coefficients. Only the following rows survive:
\[
\begin{array}{c|c|c}
 \text{family}&(A,B,C,D,r,w,z,l,m,n)&p(x)\\ \hline
 \sqrt3&(0,3,-4,8,1,-2,4,-2,3,4)&x^2-672x+2304\\
 \sqrt3&(0,3,4,8,1,2,4,2,3,4)&x^2-672x+2304\\
 \sqrt6&(3,3,4,4,2,2,2,2,3,2)&x^2-480x+2304
\end{array}
\]
Consequently $\delta=336\pm192\sqrt3$ or
$\delta=240\pm96\sqrt6$, respectively. These are the source
Frobenius--Perron dimensions and their quadratic conjugates.
Global dimension is Galois-equivariant, and the based ring does
not change under Galois conjugation \cite{ENO}. Choosing the
conjugation that sends $\delta$ to the Frobenius--Perron value
gives a pseudo-unitary category.
\end{proof}

Their full induced-unit compression quantity equals one, satisfying the
necessary bound. Saturation concerns the projection of the modular
$S$-image onto the twist-one space; it does not put the full image there.
The stronger equality constraint below, combined with two induction
trace moments, will exclude all three rings in the next section.

\subsection{A stronger consequence of compression equality}

There is nevertheless a useful restriction at equality. Let
$\mathsf S,\mathsf T$ be a unitary representation of
$SL_2(\Z/N\Z)$ and let $P$ project onto the $\mathsf T$-fixed space.
Suppose $w$ is fixed by $\mathsf T$ and has diagonal-torus character
$\chi$ of conductor $q$. Then
\[
 \|P\mathsf S w\|^2=\|w\|^2/q
 \quad\Longrightarrow\quad
 \mathsf T^q\mathsf S w=\mathsf S w.
\]
In fact the cyclic subrepresentation generated by $w$ factors through
level $q$. To see this, refine the local proof of
\cite[Lemma~4.1]{SchopierayDimensionBlock}. At a prime power $\ell^a$,
each irreducible with a unipotent-fixed vector embeds in the permutation
representation on primitive vectors. For a nontrivial character of
conductor $\ell^c$, the maximal singular subspace of the compressed
Weyl operator is spanned by the two functions denoted $F_\alpha$ and
$G_{\overline\alpha}$ there. Both depend only on the coordinates
modulo $\ell^c$. Every irreducible containing an equality vector thus
inflates from that level. For a trivial local character, equality says
that both $w$ and $\mathsf S w$ are unipotent-fixed; upper and lower
unipotents then fix $w$, so its local representation is trivial.
Orthogonal irreducible decomposition and the tensor-product description
over prime powers give the assertion. The complete local argument and
an independent audit are retained in the compression supplement.

For the three rings above, use their pseudo-unitary Galois conjugates,
put $q=12,12,24$, respectively, and let
$w=e_{\mathbf1}-e_{A_\sigma}$ be the difference of the two
dimension-character constituents of $I(\mathbf1)$. Their dimension
field is quadratic of conductor $q$, and their other formal codegrees
are rational. On the known induced-unit coordinates,
\[
 \operatorname{proj}_{I(\mathbf1)}(\mathsf S w)
 =\left(D^{-1}-\overline D^{-1}\right)w
 =-\frac{1}{\sqrt q}w.
\]
This already attains the compression bound, so the preceding argument
applies to the full twist-one projection. The coordinate at a simple
central object $V$ is
$\dim(V)/D-\sigma(\dim(V)/D)$. Consequently
\[
 \dim(V)/D\notin\mathbb Q
 \quad\Longrightarrow\quad \theta_V^q=1.
\]
This constrains twists of specific central objects; it does not assert
that every central twist has order dividing $q$. The next section keeps
that distinction when applying the two-moment obstruction.
